\respcap{cap:pd}
\resp{pd-max} \textbf{Sim.} Se a última multiplicação de uma parentização máxima de $M_i\cdots M_j$ é feita no ponto $k$, o custo é $c(i..k)+c(k+1..j)+r_{i-1}r_kr_j$, e o último termo depende só de $k$, não de como cada parte foi parentizada. Se uma das partes não fosse máxima, trocá-la por uma máxima aumentaria o total, uma contradição. A recorrência é a mesma, com $\max$ no lugar de $\min$, e o custo é $\Theta(n^3)$.

\resp{pd-lis} Seja $L[i]$ o tamanho da maior subsequência crescente que \emph{termina em} $a_i$: $L[i]=1+\max\{L[j]:j<i,\ a_j<a_i\}$, com $\max\emptyset=0$; a resposta é $\max_iL[i]$, e guardar o $j$ escolhido em $ant[i]$ permite reconstruí-la. Custo $O(n^2)$. Por exemplo, $5,2,8,6,3,6,9,7$ dá $L=1,1,2,2,2,3,4,4$ e resposta 4 ($2,3,6,9$). Em $O(n\log n)$: mantenha $d[\ell]=$ menor final possível de uma subsequência crescente de tamanho $\ell$, que é um vetor crescente; para cada $a_i$, encontre por busca binária o primeiro $d[\ell]\ge a_i$ e substitua-o por $a_i$ (ou acrescente $a_i$ no fim, se não houver). O tamanho final de $d$ é a resposta.

\resp{pd-pal} Seja $P(i,j)$ o tamanho do maior palíndromo que é subsequência de $s_i..s_j$: $P(i,i)=1$ e $P(i,i-1)=0$; se $s_i=s_j$, $P(i,j)=2+P(i+1,j-1)$; senão, $P(i,j)=\max\{P(i+1,j),P(i,j-1)\}$. Preencha por tamanho crescente do trecho: $\Theta(n^2)$ de tempo e memória. Para ``character'', o valor é 5 (``carac''). Esse valor também é o tamanho da subsequência comum máxima entre $s$ e $s$ invertida.

\resp{pd-dag} Ordene os vértices topologicamente (Capítulo~\ref{cap:dfsdig}). Faça $D(s)=0$ e $D(v)=-\infty$ para os demais; na ordem topológica, $D(v)=\max_{uv\in E}\{D(u)+w(u,v)\}$, guardando o predecessor que deu o máximo. A resposta é $D(t)$, em $\Theta(n+m)$. Num grafo acíclico todo caminho é simples, e por isso a subestrutura ótima vale; em grafos gerais, o caminho simples mais longo é um problema NP-difícil.

\resp{pd-ilim} Para $w=0..W$: \[K(w)=\max\big\{0,\ \max_{i:\,p_i\le w}\{K(w-p_i)+v_i\}\big\}.\] $\Theta(nW)$. Por exemplo, com itens $(p,v)=(3,5),(4,7),(2,3)$ e $W=10$, o máximo é 17 (por exemplo, $4+4+2$). É o mesmo que a mochila 0-1 em vetor único, com $w$ \emph{crescente}.

\resp{pd-corte} Ordene os pontos e acrescente as sentinelas $L_0=0$ e $L_{m+1}=n$. Seja $C(i,j)$ o custo mínimo de fazer todos os cortes estritamente entre $L_i$ e $L_j$: $C(i,i+1)=0$ e $C(i,j)=(L_j-L_i)+\min_{i<k<j}\{C(i,k)+C(k,j)\}$, em que $L_k$ é o primeiro corte feito no trecho. A resposta é $C(0,m+1)$; guardando o $k$ ótimo, a sequência de cortes sai em pré-ordem. Custo $\Theta(m^3)$. No exemplo, o ótimo é \textbf{38}: corte primeiro em 10 e depois em 2 e 8, em qualquer ordem (as seis ordens custam $50,46,40,40,38,38$).

\resp{pd-rec} Percorra a tabela a partir de $(n,m)$: se $x_i=y_j$, coloque $x_i$ no início da resposta e vá para $(i-1,j-1)$; senão, vá para a vizinha, $(i-1,j)$ ou $(i,j-1)$, de maior valor. Custo $O(n+m)$. No exemplo, subindo nos empates, obtém-se \textbf{BCBA}; as subsequências comuns de tamanho 4 são BCAB, BCBA e BDAB (§\ref{sec:lcs}).

\resp{pd-mon} (a) Para $r=(10,100,5,50)$: $m[1,2]=5000$, $m[2,3]=25\,000$ e $m[1,3]=\min\{75\,000,\,7500\}=7500$, com a ordem $((M_1M_2)M_3)$. Para $r=(20,50,10,1,100)$: $m[1,4]=3500$, com a ordem $((M_1(M_2M_3))M_4)$; a tabela completa está na §\ref{sec:cadeia}. (b) A última linha é $M(4,\cdot)=0,2,9,9$; o ótimo é 9, com os itens 1 e 2. (c) O tempo depende do \emph{valor} de $W$, mas $W$ ocupa só $\log_2W$ bits: o algoritmo é exponencial no tamanho da entrada.

\resp{pd-apa} (a) Tamanho 4: \textrm{acda} ou \textrm{dcda}. (b) 240, com os itens 1, 2 e 4 (pesos $20+30+50=100$). (c) Algoritmo de Kadane: $S(i)=\max\{a_i,\ S(i-1)+a_i\}$ é a melhor soma de um subvetor que \emph{termina} em $i$, e a resposta é $\max_iS(i)$, em $O(n)$. No exemplo, 55, com o subvetor $10,-5,40,10$.

\resp{pd-l6a} (a) Custo 2010, com $((A_1A_2)((A_3A_4)(A_5A_6)))$. (b) Guarde em $s[i,j]$ o $k$ que deu o mínimo e use \textsc{Multiplica}$(A,s,i,j)$ da §\ref{sec:cadeia}. (c) Na reconstrução, quando $x_i\ne y_j$ e $C[i-1,j]=C[i,j-1]$, siga os dois caminhos recursivamente e elimine as repetições; pode haver exponencialmente muitas respostas.

\resp{pd-l6b} (a) O número de Catalan $C_6=\binom{12}6/7=132$. (b) $c(i,j)=\sum_{k=i}^ja_k+\min_{i\le r\le j}\{c(i,r-1)+c(r+1,j)\}$, com $c(i,i-1)=0$, em $\Theta(n^3)$. Resultado: custo 185, com raiz \emph{if}; à esquerda, \emph{case}, com o caminho \emph{else}$\to$\emph{end} à direita; à direita da raiz, \emph{then}, com \emph{of} como filho esquerdo.

\resp{pd-l6c} (a) $S[i]=a_i+\max\{S[j]:j<i,\ a_j<a_i\}$ (0 se não houver tal $j$), e a resposta é $\max_iS[i]$; por exemplo, $1,101,2,3,100,4,5$ dá 106 ($1,2,3,100$).\\ (b) Em pós-ordem: $com(v)=s(v)+\sum_fsem(f)$ e $sem(v)=\sum_f\max\{com(f),sem(f)\}$, somando sobre os filhos $f$; a resposta é $\max\{com,sem\}$ na raiz, em $\Theta(n)$.\\ (c) É a subsequência comum máxima das duas sequências.\\ (d) $f(1)=f(2)=2$ e $f(n)=f(n-1)+f(n-2)$. A última faixa é branca ou vermelha; se a penúltima não é azul, a sequência vem de uma de tamanho $n-1$, estendida de um único jeito; se é azul, vem de uma de tamanho $n-2$, estendida pela faixa azul e pela cor que falta. Valores: $2,2,4,6,10,16,26,\dots$

\resp{pd-mc458} (a) Capacidades $0..14$:\\
\hspace*{1em}item 3 (peso 5): \texttt{O - O - - I - O - O - - I - I}\\
\hspace*{1em}item 4 (peso 2): \texttt{O - O - I O - O - O - I O - O}\\
(b) Um único vetor $K[0..W]$: para cada item, percorra $w$ \emph{decrescendo} de $W$ até $p_i$, fazendo $K[w]:=\max\{K[w],K[w-p_i]+v_i\}$. A memória cai para $O(K)$; com $w$ crescendo, o item poderia ser reaproveitado. (c) Fixe o par de linhas $(i_1,i_2)$, some cada coluna entre elas e aplique Kadane ao vetor de somas: $O(m^2n)$, usando somas de prefixo por coluna. (d) $P(a,b)=\max\{p_{ab},\ \max_{a'<a}P(a',b)+P(a-a',b),\ \max_{b'<b}P(a,b')+P(a,b-b')\}$ para todos os $a\le n$ e $b\le m$: $O(nm(n+m))\subseteq O(n^2m^2)$.

\resp{pd-edit} $D(i,0)=i$, $D(0,j)=j$ e $D(i,j)=\min\{D(i-1,j)+1,\ D(i,j-1)+1,\ D(i-1,j-1)+[x_i\ne y_j]\}$, em $\Theta(nm)$. De ``kitten'' para ``sitting'' a distância é \textbf{3}: trocar k por s, trocar e por i e inserir g.

\resp{pd-troco} São 2 moedas ($3+3$). O guloso pega $4+1+1$, com 3 moedas. Programação dinâmica: $T(0)=0$ e $T(s)=1+\min_{c\le s}T(s-c)$; os valores de $T(0)$ a $T(10)$ são $0,1,2,1,1,2,2,2,2,3,3$. Custo $\Theta(nS)$.

\resp{pd-haste} $r(0)=0$ e $r(n)=\max_{1\le i\le n}\{p_i+r(n-i)\}$. Os valores de $r(1)$ a $r(8)$ são $1,5,8,10,13,17,18,22$: o máximo para 8 é \textbf{22}, com os pedaços $2+6$.

\resp{pd-grade} São 6 passos, 3 para a direita e 3 para baixo, em qualquer ordem: $\binom63=20$ caminhos. Programação dinâmica: $P(i,j)=P(i-1,j)+P(i,j-1)$, com $P(1,1)=1$ e $P=0$ nas casas bloqueadas. Com a casa $(2,2)$ bloqueada, sobram 8 caminhos.

\resp{pd-part} Se a soma $S$ é ímpar, não é possível. Senão, decida se algum subconjunto soma $S/2$: $R(i,s)=R(i-1,s)\lor R(i-1,s-a_i)$, em $\Theta(nS)$, pseudopolinomial. $\{3,1,1,2,2,1\}$: sim ($3+2=5$ de um lado). $\{2,3,7\}$: a soma é 12, mas nenhum subconjunto soma 6: não.

\resp{pd-catalan} A última multiplicação separa os $k$ primeiros fatores dos $n-k$ últimos, e cada lado pode ser parentizado de forma independente: $P(n)=\sum_{k=1}^{n-1}P(k)P(n-k)$. Os valores para $n=1,\dots,6$ são $1,1,2,5,14,42$ --- os números de Catalan $C_{n-1}$.

\resp{pd-lcsmem} Troque $X$ e $Y$ se for preciso, para que $Y$ seja a menor. Mantenha dois vetores, $ant[0..m]$ e $atual[0..m]$, com as linhas $i-1$ e $i$ da tabela: $atual[0]:=0$ e, para $j=1..m$, $atual[j]:=1+ant[j-1]$ se $x_i=y_j$ e $\max\{ant[j],atual[j-1]\}$ caso contrário; ao fim de cada linha, troque os papéis dos vetores. A resposta é o último valor calculado. Memória $O(\min\{n,m\})$ e tempo $\Theta(nm)$.

\resp{pd-formas} Seja $F(s)$ o número de maneiras; processando as moedas uma de cada vez (para não contar ordens diferentes como maneiras diferentes): para cada moeda $c$ e para $s=c,\dots,S$, faça $F(s):=F(s)+F(s-c)$, com $F(0)=1$. Para 5 há \textbf{4} maneiras: $5$; $2+2+1$; $2+1+1+1$; $1+1+1+1+1$. Para 10 há 10 maneiras.

\resp{pd-substring} Seja $L(i,j)$ o tamanho do maior substring comum que \emph{termina} em $x_i$ e em $y_j$: $L(i,j)=L(i-1,j-1)+1$ se $x_i=y_j$ e $0$ caso contrário. A resposta é o maior $L(i,j)$, em $O(nm)$. Para ``ABABC'' e ``BABCA'', o valor é 4, com o substring ``BABC''.
